\documentclass[10pt,a4paper]{amsart}
\usepackage[margin=19mm]{geometry}
\usepackage{amsmath,amssymb,mathtools}
\usepackage[scr=rsfs,cal=euler]{mathalfa}
\usepackage{xfrac}
\usepackage[unicode,hidelinks,bookmarks=false]{hyperref}
\newcommand{\ie}{i.e.\ }
\newcommand{\cf}{cf.\ }
\newcommand{\resp}{respectively\ }
\newcommand{\Subsection}{\subsection}
\providecommand{\nobreakdash}{\nobreak\hbox{-}}
\setlength{\emergencystretch}{2em}
\multlinegap0pt
\usepackage{amssymb}
\usepackage{xcolor}
\newtheorem{lemma}{Lemma}
\newtheorem{theorem}{Theorem}
\newtheorem{remark}{Remark}
\newtheorem{corollary}{Corollary}
\newcommand{\C}{\mathbb{C}}
\newcommand{\R}{\mathbb{R}}
\title{Elementary Quantum Gates from Lie Group Embeddings in \texorpdfstring{$U(2^n)$}{U(2^n)}: Geometry, Universality, and Discretization}
\author{Antonio Falco, Daniela Falco-Pomares, Hermann G. Matthies}
\date{Source: arXiv:2601.17936v2 (CC BY 4.0)}
\begin{document}
\maketitle

\noindent\emph{Source and licence.} Elementary Quantum Gates from Lie Group Embeddings in \texorpdfstring{$U(2^n)$}{U(2^n)}: Geometry, Universality, and Discretization. Version arXiv:2601.17936v2 (CC BY 4.0), distributed by arXiv under the Creative Commons Attribution 4.0 International licence, http://creativecommons.org/licenses/by/4.0/, as recorded in the arXiv licence metadata for this exact version and shown on the abstract page https://arxiv.org/abs/2601.17936v2. Full licence notice: \texttt{licenses/arxiv-2601.17936-CC-BY-4.0.txt}.\par\smallskip

\noindent\textit{Editorial scope.} Counted unit: the article's theorem thm:UN\_compilation\_SK (source0.tex lines 1304--1328). The article's complete original proof of it is delivered with it (source0.tex lines 1330--1376, one \texttt{proof} environment of 47 lines). No mathematics is written, repaired or re-derived: the statement and the proof are the article's own lines, in the article's own notation, and no retained line is substituted or re-flowed.  Retained uncounted as the source-local context the proof needs: lines 993--1007; lines 1034--1044; lines 1037--1041; lines 1196--1200; lines 1209--1220; lines 1245--1253; lines 1273--1281; lines 1378--1464.  Nothing was omitted from any retained span, so the package carries no omission record.  No retained line contains a citation, so the wrapper carries no bibliography and no bibliography entry.  No retained line contains a figure environment, an image-inclusion command or drawing code, so no image is delivered and the package declares no asset dependency.  Editorial formatting only, in addition: an AMS \texttt{amsart} preamble that restates the article's own declarations and macros used by the retained lines, namely the environments \texttt{lemma}, \texttt{theorem}, \texttt{remark}, \texttt{corollary} on separate counters, together with the standard packages those lines need.  The retained environments print in this document as Lemma 1, Theorem 1, Remark 1, Corollary 1 in order of appearance, whereas the article prints the same items with the numbering of its own full document.  The complete native source of the article as distributed in the arXiv e-print for this exact version is bundled unchanged as \texttt{sources/193352/source0.tex}.
\begin{lemma}[The special-unitary torus is generated by $SU(2)$ two-level phases]
\label{lem:T0_generated_by_two_level_phases_concise}
For each $j=2,\dots,N$ define
\[
H_{1j}:=\frac{i}{2}\big(E_{11}-E_{jj}\big)\in\mathfrak{su}(N),
\qquad
\gamma_{1j}(t):=\exp(tH_{1j})\in SU(N).
\]
Then $\gamma_{1j}(t)\in SU(N)[W_{1j}]$ for $W_{1j}=\mathrm{span}\{e_1,e_j\}$, and
\[
T_0=\Big\langle\,\gamma_{12}(\R),\dots,\gamma_{1N}(\R)\,\Big\rangle.
\]
Consequently, every $D\in T$ admits a factorization $D=e^{i\theta}D_0$ with $\theta\in\R$ and
$D_0\in\big\langle\,SU(N)[W_{12}]\allowbreak\cup\cdots\allowbreak\cup SU(N)[W_{1N}]\,\big\rangle$.
\end{lemma}

\begin{theorem}[Unitary QR/Givens factorization into coordinate two-level gates]
\label{thm:two_level_factorization_concise}
Fix an orthonormal basis $\{e_1,\dots,e_N\}$ of $\C^N$. Then every $U\in U(N)$ admits a factorization
\begin{equation}\label{eq:U_two_level_factorization_concise}
U=\Big(\prod_{k=1}^{K} T_k\Big)\,D,
\qquad
T_k\in U(N)[W_{p_k,q_k}],\ \ D\in T,
\end{equation}
where each $W_{p_k,q_k}=\mathrm{span}\{e_{p_k},e_{q_k}\}$ is a coordinate two-plane and $K\le \frac{N(N-1)}{2}$.
Moreover, the factors can be chosen by standard Givens elimination to annihilate prescribed subdiagonal entries.
\end{theorem}

\begin{equation}
U=\Big(\prod_{k=1}^{K} T_k\Big)\,D,
\qquad
T_k\in U(N)[W_{p_k,q_k}],\ \ D\in T,
\end{equation}

\begin{remark}[Determinants and diagonal phases]\label{rem:det_bookkeeping}
Each $V\in U(2)$ admits a factorization $V=e^{i\theta}\widetilde V$ with $\widetilde V\in SU(2)$ (take $e^{i\theta}=\sqrt{\det V}$).
In the exact synthesis \eqref{eq:U_two_level_factorization_concise}, the scalar phases $e^{i\theta_k}$ may be absorbed into the trailing diagonal
unitary $D$. Accordingly, the approximation results below are stated for $SU(2)$ factors and extend to $U(2)$ with this bookkeeping.
\end{remark}

\begin{theorem}[Solovay--Kitaev theorem for $SU(2)$]\label{thm:SK}
Let $G_2\subset SU(2)$ be a finite set such that the subgroup $\langle G_2\rangle$ is dense in $SU(2)$.
Then there exist constants $c,C>0$, depending only on $G_2$, and a deterministic classical algorithm with the following property:
for every $V\in SU(2)$ and every $\varepsilon\in(0,1)$ the algorithm produces a word $w=w(V,\varepsilon)$ over
$G_2\cup G_2^{-1}$ such that
\[
\|V-w\|\le \varepsilon,
\qquad
|w|\le C\,\log^{c}\!\Big(\frac{1}{\varepsilon}\Big).
\]
Moreover, the running time is polynomial in $\log(1/\varepsilon)$.
\end{theorem}

\begin{lemma}[Operator-norm isometry for coordinate two-level embeddings]
\label{lem:two_level_isometry}
Let $1\le p<q\le N$ and let $\phi_{p,q}:U(2)\hookrightarrow U(N)$ be the coordinate two-level embedding supported on
$W_{p,q}=\mathrm{span}\{e_p,e_q\}$. Then for all $V,W\in U(2)$,
\[
\big\|\phi_{p,q}(V)-\phi_{p,q}(W)\big\|=\|V-W\|.
\]
In particular, $\phi_{p,q}$ is an isometric embedding with respect to the operator norm.
\end{lemma}

\begin{lemma}[Telescoping bound for products]
\label{lem:error_accumulation}
Let $A_1,\dots,A_K,B_1,\dots,B_K\in U(N)$. Then
\[
\Big\|\prod_{k=1}^K A_k-\prod_{k=1}^K B_k\Big\|
\ \le\
\sum_{k=1}^K \|A_k-B_k\|.
\]
\end{lemma}

\begin{remark}[Worst-case scaling in $n$]\label{rem:worst_case_scaling}
Since $K\le N(N-1)/2=\Theta(N^2)=\Theta(4^n)$, the bound \eqref{eq:length_bound_global} yields worst-case word lengths that scale
exponentially in $n$ for generic targets in $U(2^n)$. This reflects the parameter complexity of $U(2^n)$ and should be interpreted
as providing a representation-invariant and modular discretization mechanism rather than an efficiency guarantee for arbitrary
instances.
\end{remark}

%============================================================
\subsection{Compiling the diagonal factor}
\label{subsec:compile_diagonal_optional}
%============================================================

In the QR/Givens synthesis of Theorem~\ref{thm:UN_compilation_SK} the target unitary $U\in U(N)$ is produced up to a trailing diagonal
unitary $D$. If one wishes to obtain a ``pure word'' over embedded $SU(2)$ generators (up to an unavoidable global phase), then $D$ can
itself be synthesized by two-level phase rotations supported on coordinate planes. We record a convenient formulation.


\begin{corollary}[Pure-word compilation up to global phase]
\label{cor:pure_word_compilation}
Under the assumptions of Theorem~\ref{thm:UN_compilation_SK}, one may absorb the diagonal factor into the word up to a global phase:
there exist $\theta\in\mathbb{R}$ and a word $\widetilde W(U,\varepsilon)$ over
$\mathcal{A}_{\mathrm{elem}}(n;G_2)\cup\mathcal{A}_{\mathrm{elem}}(n;G_2)^{-1}$ such that
\[
\big\|U-e^{i\theta}\widetilde W(U,\varepsilon)\big\|\le \varepsilon.
\]
\end{corollary}

\begin{proof}
Apply Theorem~\ref{thm:UN_compilation_SK} to obtain $U=W(U,\varepsilon)\,D+E$ with $\|E\|\le\varepsilon$.
Decompose $D=e^{i\theta}D_0$ as in Lemma~\ref{lem:T0_generated_by_two_level_phases_concise} and write $D_0$ as a finite product of
elements from $\phi_{1j}(SU(2))$ (hence as a word over the corresponding embedded alphabet).
Let $\widetilde W(U,\varepsilon):=W(U,\varepsilon)\cdot(\text{word for }D_0)$. Then $e^{i\theta}\widetilde W=W D$ by construction, so
$U-e^{i\theta}\widetilde W=E$ and the norm bound follows.
\end{proof}

%============================================================
\subsection{Algorithmic summary and interpretation}
\label{subsec:algorithmic_summary}
%============================================================

We summarize the construction of Theorem~\ref{thm:UN_compilation_SK} as an explicit compilation procedure.

\medskip
\noindent\textbf{Compilation procedure (coordinate two-level factorization + Solovay--Kitaev).}
Given $U\in U(N)$, $\varepsilon\in(0,1)$, and a finite universal set $G_2\subset SU(2)$, the following steps produce
a word $W$ over $\mathcal{A}_{\mathrm{elem}}(n;G_2)\cup\mathcal{A}_{\mathrm{elem}}(n;G_2)^{-1}$ and a diagonal $D\in U(N)$ such that
$\|U-WD\|\le \varepsilon$.

\begin{quote}
\begin{enumerate}
\item \emph{Exact two-level factorization.}
Compute a coordinate Givens/QR factorization
\[
U=\Big(\prod_{k=1}^{K}\phi_{p_k,q_k}(V_k)\Big)\,D,
\qquad 1\le p_k<q_k\le N,\ \ V_k\in U(2),
\]
with $K\le N(N-1)/2$.
If desired, write $V_k=e^{i\theta_k}\widetilde V_k$ with $\widetilde V_k\in SU(2)$ and absorb the phases $e^{i\theta_k}$ into $D$.

\item \emph{Local approximation in $SU(2)$.}
Set $\delta:=\varepsilon/K$.
For each $k=1,\dots,K$, apply a Solovay--Kitaev routine to obtain a word $w_k$ over $G_2\cup G_2^{-1}$ such that
\[
\|\,\widetilde V_k-w_k\,\|\le \delta.
\]

\item \emph{Lift and concatenate.}
Define
\[
W:=\prod_{k=1}^{K}\phi_{p_k,q_k}(w_k),
\]
where $\phi_{p_k,q_k}(w_k)$ denotes the word obtained by replacing each letter of $w_k$ by its embedded generator in
$\phi_{p_k,q_k}(G_2)\cup \phi_{p_k,q_k}(G_2)^{-1}$.
\end{enumerate}
\end{quote}

\noindent
By Lemmas~\ref{lem:two_level_isometry} and~\ref{lem:error_accumulation}, the resulting pair $(W,D)$ satisfies $\|U-WD\|\le\varepsilon$.
\medskip

\begin{remark}[Interpretation in the geometric framework]
The factorization step may be viewed as selecting a finite sequence of logical two-planes (here, the coordinate planes
$W_{p_k,q_k}$) and prescribing on each plane a continuous $U(2)$ motion $V_k$. The Solovay--Kitaev step then replaces each
continuous $SU(2)$ motion by a short word over a fixed finite alphabet, and the lifting step transports these words to $U(N)$ through
the corresponding two-level embeddings. In this sense, the procedure separates a \emph{continuous} choice of support (a path through
logical qubits) from a \emph{discrete} approximation of the induced $SU(2)$ dynamics.
\end{remark}

\begin{theorem}[Compilation via two-level factorization and Solovay--Kitaev]
\label{thm:UN_compilation_SK}
Fix $N=2^n$ and let $G_2\subset SU(2)$ be a finite universal set.
Let $U\in U(N)$ and $\varepsilon\in(0,1)$. Then there exist:
\begin{itemize}
\item a word $W(U,\varepsilon)$ over the finite alphabet
\[
\mathcal{A}_{\mathrm{elem}}(n;G_2)\ \cup\ \mathcal{A}_{\mathrm{elem}}(n;G_2)^{-1},
\]
\item and a diagonal unitary $D\in U(N)$,
\end{itemize}
such that
\begin{equation}\label{eq:global_approx}
\big\|U - W(U,\varepsilon)\,D\big\|\ \le\ \varepsilon.
\end{equation}
Moreover, one may choose $W(U,\varepsilon)$ with length bounded by
\begin{equation}\label{eq:length_bound_global}
|W(U,\varepsilon)|\ \le\ K\cdot C\,\log^{\,c}\!\Big(\frac{K}{\varepsilon}\Big),
\qquad
K\le \frac{N(N-1)}{2},
\end{equation}
where $c,C$ are the Solovay--Kitaev constants from Theorem~\ref{thm:SK}.
If one wishes to eliminate the residual diagonal factor, it can be compiled (up to global phase) using two-level phase rotations as in
Section~\ref{subsec:compile_diagonal_optional}.
\end{theorem}

\begin{proof}
By Theorem~\ref{thm:two_level_factorization_concise} there exist indices $(p_k,q_k)$, factors $V_k\in U(2)$, and a diagonal unitary $D\in U(N)$
such that
\[
U=\Big(\prod_{k=1}^K \phi_{p_k,q_k}(V_k)\Big)\,D,
\qquad
K\le \frac{N(N-1)}{2}.
\]
Using Remark~\ref{rem:det_bookkeeping}, write $V_k=e^{i\theta_k}\widetilde V_k$ with $\widetilde V_k\in SU(2)$ and absorb the phases
into $D$, so that we may assume $V_k\in SU(2)$ without loss of generality.

Set $\delta:=\varepsilon/K$. For each $k$, Theorem~\ref{thm:SK} yields a word $w_k$ over $G_2\cup G_2^{-1}$ such that
\[
\|V_k-w_k\|\le \delta,
\qquad
|w_k|\le C\,\log^{c}\!\Big(\frac{1}{\delta}\Big)=C\,\log^{c}\!\Big(\frac{K}{\varepsilon}\Big).
\]
Define the lifted word in $U(N)$ by
\[
W(U,\varepsilon):=\prod_{k=1}^K \phi_{p_k,q_k}(w_k),
\]
where $\phi_{p_k,q_k}(w_k)$ denotes the word obtained by embedding each letter of $w_k$ through $\phi_{p_k,q_k}$.
By Lemma~\ref{lem:two_level_isometry},
\[
\big\|\phi_{p_k,q_k}(V_k)-\phi_{p_k,q_k}(w_k)\big\|
=
\|V_k-w_k\|
\le \delta.
\]
Applying Lemma~\ref{lem:error_accumulation} gives
\[
\Big\|\prod_{k=1}^K\phi_{p_k,q_k}(V_k)-\prod_{k=1}^K\phi_{p_k,q_k}(w_k)\Big\|
\le
\sum_{k=1}^K \delta
=
\varepsilon.
\]
Right-multiplication by $D$ preserves the operator norm (unitary invariance), hence \eqref{eq:global_approx} follows.
Finally,
\[
|W(U,\varepsilon)|
=\sum_{k=1}^K |w_k|
\le
K\cdot C\,\log^{c}\!\Big(\frac{K}{\varepsilon}\Big),
\]
which is \eqref{eq:length_bound_global}.
\end{proof}

\end{document}
